% Options for packages loaded elsewhere
\PassOptionsToPackage{unicode}{hyperref}
\PassOptionsToPackage{hyphens}{url}
%
\documentclass[
  10pt,
]{article}
\usepackage{amsmath,amssymb}
\usepackage{setspace}
\usepackage{iftex}
\ifPDFTeX
  \usepackage[T1]{fontenc}
  \usepackage[utf8]{inputenc}
  \usepackage{textcomp} % provide euro and other symbols
\else % if luatex or xetex
  \usepackage{unicode-math} % this also loads fontspec
  \defaultfontfeatures{Scale=MatchLowercase}
  \defaultfontfeatures[\rmfamily]{Ligatures=TeX,Scale=1}
\fi
\usepackage{lmodern}
\ifPDFTeX\else
  % xetex/luatex font selection
  \setmainfont[]{DejaVu Serif}
  \setmonofont[]{DejaVu Sans Mono}
\fi
% Use upquote if available, for straight quotes in verbatim environments
\IfFileExists{upquote.sty}{\usepackage{upquote}}{}
\IfFileExists{microtype.sty}{% use microtype if available
  \usepackage[]{microtype}
  \UseMicrotypeSet[protrusion]{basicmath} % disable protrusion for tt fonts
}{}
\makeatletter
\@ifundefined{KOMAClassName}{% if non-KOMA class
  \IfFileExists{parskip.sty}{%
    \usepackage{parskip}
  }{% else
    \setlength{\parindent}{0pt}
    \setlength{\parskip}{6pt plus 2pt minus 1pt}}
}{% if KOMA class
  \KOMAoptions{parskip=half}}
\makeatother
\usepackage{xcolor}
\usepackage[top=.5in,bottom=.5in,left=.2in,right=.2in]{geometry}
\usepackage{longtable,booktabs,array}
\usepackage{calc} % for calculating minipage widths
% Correct order of tables after \paragraph or \subparagraph
\usepackage{etoolbox}
\makeatletter
\patchcmd\longtable{\par}{\if@noskipsec\mbox{}\fi\par}{}{}
\makeatother
% Allow footnotes in longtable head/foot
\IfFileExists{footnotehyper.sty}{\usepackage{footnotehyper}}{\usepackage{footnote}}
\makesavenoteenv{longtable}
\setlength{\emergencystretch}{3em} % prevent overfull lines
\providecommand{\tightlist}{%
  \setlength{\itemsep}{0pt}\setlength{\parskip}{0pt}}
\setcounter{secnumdepth}{-\maxdimen} % remove section numbering
\ifLuaTeX
  \usepackage{selnolig}  % disable illegal ligatures
\fi
\usepackage{bookmark}
\IfFileExists{xurl.sty}{\usepackage{xurl}}{} % add URL line breaks if available
\urlstyle{same}
\hypersetup{
  hidelinks,
  pdfcreator={LaTeX via pandoc}}

\author{}
\date{}

\begin{document}

\setstretch{1}
\section{System Functionality and Documentation (60
points)}\label{system-functionality-and-documentation-60-points}

See attached zip file. The directory structure is as follows:

\begin{verbatim}

├── sim
├── sim_pipeline_three
├── sim_pipelined
├── synth
├── synth_pipeline_three
└── synth_pipelined
\end{verbatim}

Each sim directory has a Makefile that can be used to simulate the
design using \texttt{Make\ run}. Debug info can be increased by
uncommented the \texttt{SetLogEnable} calls in \texttt{Cordic.vhd} and
\texttt{Cordic\_TB.vhd} in each directory. The synth directories contain
mostly the same files as their corresponding sim directories, but with a
few modifications because Xilinx ISE 14.7 does not support the full
VHDL-2008 standard. The printout contains only the files from the
\texttt{synth\_pipeline\_three} directory. The rest of the directories
are also not as fully commented as this.

\section{System Quality factor (25
points)}\label{system-quality-factor-25-points}

\begin{itemize}
\tightlist
\item
  Pipelined quality factor:
  \(10\frac{f'_{\text{max}}}{f_{\text{max}}}-l^{2}\) where
  \(f'_{\text{max}}\) is the max frequency of the pipelined system,
  \(f_{\text{max}}\) is the max frequency of the original system, and
  \(l^2\) is the latency in clocks.
\end{itemize}

\subsection{Stats:}\label{stats}

\begin{longtable}[]{@{}
  >{\raggedright\arraybackslash}p{(\columnwidth - 14\tabcolsep) * \real{0.1486}}
  >{\raggedright\arraybackslash}p{(\columnwidth - 14\tabcolsep) * \real{0.1419}}
  >{\raggedright\arraybackslash}p{(\columnwidth - 14\tabcolsep) * \real{0.1486}}
  >{\raggedright\arraybackslash}p{(\columnwidth - 14\tabcolsep) * \real{0.1486}}
  >{\raggedright\arraybackslash}p{(\columnwidth - 14\tabcolsep) * \real{0.0811}}
  >{\raggedright\arraybackslash}p{(\columnwidth - 14\tabcolsep) * \real{0.0743}}
  >{\raggedright\arraybackslash}p{(\columnwidth - 14\tabcolsep) * \real{0.1486}}
  >{\raggedright\arraybackslash}p{(\columnwidth - 14\tabcolsep) * \real{0.1081}}@{}}
\toprule\noalign{}
\begin{minipage}[b]{\linewidth}\raggedright
Pipelining Stages
\end{minipage} & \begin{minipage}[b]{\linewidth}\raggedright
Size (slices)
\end{minipage} & \begin{minipage}[b]{\linewidth}\raggedright
Slice Registers Used
\end{minipage} & \begin{minipage}[b]{\linewidth}\raggedright
Minimum Clock Speed
\end{minipage} & \begin{minipage}[b]{\linewidth}\raggedright
Frequency
\end{minipage} & \begin{minipage}[b]{\linewidth}\raggedright
Latency
\end{minipage} & \begin{minipage}[b]{\linewidth}\raggedright
Throughput
\end{minipage} & \begin{minipage}[b]{\linewidth}\raggedright
Quality Factor
\end{minipage} \\
\midrule\noalign{}
\endhead
\bottomrule\noalign{}
\endlastfoot
None & 3,917 / 8,672 (44\%) & 91 / 17,344 (1\%) & 172/175 nS & 5.808 MHz
& 1 clock & 5.808 million ops/s & N/A \\
1 & 2,531 / 8,672 (29\%) & 158 / 17,344 (1\%) & 89.022 nS & 11.233 MHz &
2 clocks & 11.233 million ops/s & 15.340 \\
3 & 2,173 / 8,672 (25\%) & 304 / 17,344 (1\%) & 45.681 nS & 21.890 & 4
clocks & 21.890 million ops/s & 21.689 \\
\end{longtable}

\begin{itemize}
\tightlist
\item
  All timing values are from post-PAR (place-and-route) reports.
\end{itemize}

\section{System Size (15 points)}\label{system-size-15-points}

Surprisingly, the pipelined designs used less slices than the
non-pipelined designs. I suspect that inserting DFFs in between stages
somehow allowed the compiler/synthesizer to reuse large amounts of
logic. This would make sense because a lot of the signals (such as the
internal signals) are generated only for a small part of the cordic
stages and then stored in DFFs. For example, in the three-stage
pipelined Cordic design, the muxxes associated with generating the
internal signals only need to generate the signals for 4 cordic stages
and then save them in DFFs for the next stage of the pipeline, rather
than make generate them for all 16 cordic stages. We do see an increase
in slice DFFs used, with each pipelining stage adding about 60 - 70
DFFs. This is expected because we need to latch 3 * 22 bits for the
outputs/inputs of each stage of pipelining we add, along with the select
lines/internal signals. For each stage of pipelining we add, we need to
latch:

\begin{itemize}
\tightlist
\item
  XInputBus(k) (22 bits)
\item
  YInputBus(k) (22 bits)
\item
  ZInputBus(k) (22 bits)
\item
  Func (5 bits)
\item
  CordSystem (2 bits)
\item
  Mode (1 bit)
\end{itemize}

This total to 74 DFFs per pipelining stage, but the actual increase is
less than that, likely due to compiler/routing optimizations and
trimming.

\end{document}
